\section{Introduction}
\label{sec:introduction}

\subsection{Le projet LapZero-Chess}

LapZero-Chess est un moteur d'\'echecs de type AlphaZero, d\'evelopp\'e pour
tenir sur une seule machine personnelle. L'architecture suit le sch\'ema
classique du papier de DeepMind \cite{alpha0316} : un r\'eseau r\'esiduel
pr\'edit une politique de coups et une valeur de position, et une recherche
arborescente de Monte-Carlo (MCTS) combine ces pr\'edictions avec des
statistiques de visite pour produire \`a la fois les coups jou\'es et les
cibles d'entra\^inement.

Le pipeline comporte trois \'etages. D'abord un pr\'e-entra\^inement supervis\'e
sur des parties de ma\^itres et de joueurs Lichess, qui donne au r\'eseau les
rudiments tactiques et positionnels. Ensuite une boucle de self-play : le
moteur joue contre lui-m\^eme, chaque partie produit des couples
(\'etat, politique de visites, r\'esultat) qui servent \`a r\'eentra\^iner le
r\'eseau. Enfin un bot UCI, d\'eploy\'e sur Lichess, qui utilise le m\^eme
moteur de recherche en partie r\'eelle.

Le c\oe{}ur du moteur est \'ecrit en C++17 et expos\'e \`a Python par pybind11.
Le r\'eseau est ex\'ecut\'e par ONNX Runtime, sur GPU CUDA quand il est
disponible, avec un repli CPU. Le MCTS, la g\'en\'eration de coups, la table de
transposition et le code de self-play vivent tous dans ce c\oe{}ur C++, ce qui
permet d'atteindre environ 1,8 million de n\oe{}uds par seconde en perft
mono-thread, et de garder la boucle d'entra\^inement en Python l\'eger.

\subsection{Le probl\`eme : le GPU attend le CPU}

La recherche MCTS \'evalue le r\'eseau sur des positions feuilles. Jusqu'en
septembre 2026, chaque simulation d\'eclenchait une inf\'erence sur une seule
position. La cons\'equence a \'et\'e mesur\'ee avec pr\'ecision : le d\'ebit
\'etait de 286 \`a 299 simulations par seconde, ind\'ependamment de la
position, du nombre de coups l\'egaux (14 en finale, 48 en milieu de jeu) et
de la profondeur de l'arbre. Une telle insensibilit\'e \`a la forme du travail
d\'enonce un co\^ut domin\'e par une latence fixe : environ 3,4\,ms pour lancer
un noyau GPU sur une position.

Le premier chantier, le batching par virtual loss, a group\'e jusqu'\`a huit
feuilles dans une seule inf\'erence. Le d\'ebit est mont\'e jusqu'\`a environ
850 simulations par seconde en ouverture, et un lot de 8 a \'et\'e valid\'e
comme le meilleur compromis qualit\'e-d\'ebit. Mais une fois ce travail fait,
le profil du temps a chang\'e de nature : l'inf\'erence elle-m\^eme repr\'esente
d\'esormais 96 \`a 99\,\% du temps total, et le reste est domin\'e par une
\'etape que le batching n'avait pas parall\'elis\'ee, la pr\'eparation s\'erielle
des descentes.

Concr\`etement, la boucle batch\'ee est la suivante : pendant qu'un unique
c\oe{}ur descend l'arbre huit fois de suite pour collecter huit feuilles, le
GPU ne fait rien ; puis, pendant que le GPU \'evalue le lot, ce c\oe{}ur attend.
Chaque vague paie donc deux attentes altern\'ees. Un indice suppl\'ementaire
est le remplissage : les lots contiennent souvent 5 \`a 7 positions utiles au
lieu de 8, parce qu'une collision ou une fin d'arbre ferme le lot en cours de
route. Chaque appel GPU partiellement rempli co\^ute sa latence fixe pour moins
de travail.

\subsection{La question pos\'ee}

Ce rapport r\'epond \`a une question d'ing\'enierie pr\'ecise : peut-on faire
pr\'eparer les descentes par plusieurs c\oe{}urs CPU sur un arbre partag\'e,
sans corruption silencieuse, sans casser le self-play, et pour quel gain
r\'eel de d\'ebit et de qualit\'e ?

La difficult\'e n'est pas d'\'ecrire une boucle parall\`ele, mais de rendre
correcte la cohabitation de plusieurs descentes dans une m\^eme structure de
donn\'ees qui n'a jamais \'et\'e con\c{c}ue pour cela. L'arbre est mut\'e en
permanence (visites, valeurs, enfants), la table de transposition est un
vecteur partag\'e, et le plateau d'\'echecs lui-m\^eme se modifie pendant la
g\'en\'eration des coups. Chaque acc\`es concurrent doit \^etre soit rendu
invisible, soit prot\'eg\'e par un protocole explicite. C'est ce travail
d'architecture, plus que le gain brut, qui constitue le sujet de fond de ce
document.

\subsection{Contributions}

Le chantier a produit quatre choses que ce rapport documente en d\'etail.

\begin{enumerate}[leftmargin=1.4em]
  \item \textbf{Une architecture par vagues} : les workers collectent en
  parall\`ele sans jamais \'ecrire les statistiques de l'arbre, puis le
  coordinateur traite le lot entier s\'equentiellement. La correction est
  obtenue par s\'eparation stricte des phases plut\`ot que par verrouillage
  g\'en\'eralis\'e.
  \item \textbf{Cinq composants r\'eutilisables et testables isol\'ement} :
  un \'evaluateur injectable, une table de transposition par snapshots, une
  garde de r\'eservation RAII sur des \'etats de n\oe{}ud atomiques, un
  ex\'ecuteur de workers persistant, et un collecteur de feuilles born\'e.
  \item \textbf{Une instrumentation et un harnais de mesure} capables de
  s\'eparer les temps par phase, de comparer des campagnes appari\'ees, et de
  d\'etecter une r\'egression de qualit\'e \`a $0,1$ point pr\`es sur
  2\,500 puzzles.
  \item \textbf{Un verdict mesur\'e} : le multic\oe{}ur apporte $+8$ \`a
  $+26\,\%$ selon la position, plafonne d\`es deux workers, ne d\'egrade pas la
  queue de latence ni la qualit\'e, et a \'et\'e activ\'e dans le bot UCI \`a
  huit workers.
\end{enumerate}

\subsection{Ce que ce document ne couvre pas}

Trois sujets voisins sont volontairement laiss\'es de c\^ot\'e. La
\emph{production continue}, qui ferait descendre les workers pendant les
inf\'erences et les remont\'ees, est l'\'etape~2 de la conception ; elle n'est
pas impl\'ement\'ee et les mesures de ce rapport servent pr\'ecis\'ement \`a
d\'ecider si elle vaut son co\^ut. L'\emph{assainissement de l'appel \`a
l'\'evaluateur} (forme de lot fixe, buffers r\'eutilis\'es, quantification) est
un chantier ind\'ependant, document\'e dans
\code{2026-09-19-cout-calcul-cpu-gpu.md}. Enfin la parall\'elisation du
self-play n'est pas concern\'ee : il parall\'elise d\'ej\`a des parties
ind\'ependantes et n'a pas besoin d'un arbre partag\'e.

\subsection{Conventions de lecture}

Les mesures de d\'ebit sont exprim\'ees en \emph{simulations par seconde}
(sims/s). Une simulation est une remont\'ee de valeur compl\`ete, qu'elle ait
co\^ut\'e une inf\'erence ou non (une simulation qui s'arr\^ete sur un n\oe{}ud
terminal ou sur un succ\`es de table n'appelle pas le r\'eseau). Les trois
positions de r\'ef\'erence du banc sont appel\'ees \emph{ouverture},
\emph{milieu} et \emph{finale} ; leurs FEN sont donn\'ees en annexe~C. Le
\emph{remplissage} est le nombre de positions r\'eellement \'evalu\'ees par
appel GPU. La \emph{table de transposition} est not\'ee TT, et la politique
\code{h0} d\'esigne la cl\'e s\'emantique du cache d\'ecrite au
chapitre~2. Les sessions de mesure varient de 15 \`a 30\,\% entre elles
(horloge et thermique du GPU) : sauf mention contraire, seules les
comparaisons intra-session sont utilis\'ees, et les protocoles sont con\c{c}us
pour annuler la d\'erive.
